\section{Introduction}
\label{sec:intro}

Event cameras report per-pixel logarithmic brightness changes with
microsecond latency, high dynamic range, and sparse output, and have
transformed visible-light high-speed and low-light
vision~\cite{rebecq2018esim,gehrig2020vid2e}. The same advantages are
compelling in the thermal infrared, where \emph{event-based} sensing would
additionally remove the frame-rate bottleneck imposed by microbolometer
time constants. Event-based infrared hardware is now arriving: the DARPA
FENCE programme targets event-based infrared focal plane arrays with
embedded neuromorphic processing at $<$1.5\,W~\cite{fence}; Raytheon
demonstrated the first event-based MWIR camera in April
2026~\cite{raytheon2026}; and uncooled Poisson-bolometer LWIR pixels with
nanosecond switching have been reported~\cite{mousa2026poisson,neuralbolometers2026}.

Yet \emph{no public thermal event dataset---real or synthetic---exists}. Our
novelty audit (Sec.~\ref{sec:related}) confirms that every nearby candidate
is either visible-light-only (EVIF~\cite{geng2024evif},
DVS-PedX~\cite{sakhai2026dvspedx}, SEVD~\cite{sevd}) or hardware
demonstrations without data release~\cite{raytheon2026,mousa2026poisson}.
This creates an algorithm vacuum precisely when hardware is about to
arrive: detection and tracking methods that will run on thermal event
cameras cannot be developed or benchmarked today.

We close this gap with a \emph{pre-hardware algorithm study}. Following the
``recycle annotated video'' template established for visible light by
vid2e~\cite{gehrig2020vid2e} and v2e~\cite{hu2021v2e}, and recently
mainstreamed by DVS-PedX~\cite{sakhai2026dvspedx}, we convert public
thermal-infrared video into synthetic thermal event streams. Thermal video,
however, is not visible-light video in a different colour map: it carries
AGC flicker, microbolometer temporal lag, NETD-limited contrast, and
radiometric quantization. Our central observation is that
\emph{thermal-specific preprocessing and sensor-physics parameterization are
first-class contributions}, not implementation details.

\paragraph{Contributions}
\begin{itemize}
  \item[\textbf{C1}] The first open-source thermal-infrared\,$\to$\,event
  conversion pipeline, with pivot-centred AGC de-flickering that provably
  recovers the radiometric event floor and a
  threshold$\leftrightarrow$NETD calibration.
  \item[\textbf{C2}] \textbf{STED}, a synthetic thermal event detection
  benchmark with ground-truth physics parameters, and \textbf{FLIR-STED},
  the first event-stream conversion of continuous FLIR ADAS video with
  COCO labels---together the first public thermal event detection data.
  \item[\textbf{C3}] A fidelity study across simulator modes and
  parameters, including an interpolation audit showing that flow-guided
  interpolation does \emph{not} transfer to thermal video, and a
  sensitivity band analysis rather than single-point conclusions.
  \item[\textbf{C4}] Three-channel detection baselines (thermal frame,
  thermal event, early fusion) with scene slicing (day/night, motion
  magnitude, target size, contrast), quantifying where the event channel
  helps.
\end{itemize}

All conclusions are stated as \emph{simulator-conditional}: we do not claim
that synthetic results equal real sensor performance, and we explicitly
bound the sim-to-real gap following the lessons
of~\cite{stoffregen2020sim2real}. Code, configs, and the FLIR-STED event
streams are released; source video remains subject to its original licences.
